\documentclass[a4paper,14pt]{article}

% текст ВКР рекомендуется писать в кодировке utf-8
\usepackage[utf8]{inputenc}

% Подключение пакета с шаблоном ВКР
\usepackage{spbpu}

% Используемые функции
% !TEX root = main.tex

% Вставка рисунка с рамкой
\newcommand{\myPicture}[4]{
	\begin{figure}[H]
		\begin{center}
			\fcolorbox{black}{white}{\includegraphics[width=#1]{#2}}
			\captionsetup{justification=centering}
			\caption{#3}
			\label{#4}
		\end{center}
	\end{figure}
}

% Вставка рисунка без рамки
\newcommand{\myPictureNoBox}[4]{
	\begin{figure}[H]
		\begin{center}
			\includegraphics[width=#1]{#2}
			\captionsetup{justification=centering}
			\caption{#3}
			\label{#4}
		\end{center}
	\end{figure}
}

%%%% Единицы измерения

% (м2·°С)/Вт
\newcommand\SIteploper{(\text{м}^2 \cdot \textdegree C)/\text{Вт}}
% Вт/(м2·°С)
\newcommand\SIteplootd{\text{Вт}/(\text{м}^2 \cdot \textdegree C)}

%%% Обозначения таблиц и руснков
%\newcommand\piclabel{Рисунок\thinspace}
\newcommand\piclabel{Рис.}
\newcommand\tablabel{Таблица\thinspace}
%\newcommand\tablabel{Табл.}
\newcommand\appendixLabel{Приложение\thinspace}




% Константы
\university{Министерство науки и высшего образования Российской Федерации Санкт-Петербургский политехнический университет Петра Великого}
\institute{Физико-механический институт}
\school{Высшая школа прикладной математики и вычислительной физики}
\thesisTitle{ВЫПУСКНАЯ КВАЛИФИКАЦИОННАЯ РАБОТА}
\title{Название работы}
\studyField{01.03.02 <<Прикладная математика и информатика>>}
\profile{01.03.02\_02 <<Системное программирование>>}
\group{5030102/10201}
\author{И. О. Фамилия}
\supervisorTitle{доцент, к.т.н.}
\supervisor{И. О. Фамилия}
\city{Санкт-Петербург}
\publishYear{2025}

\begin{document}
	% 1. Титульный лист
    \thispagestyle{empty} % убираем колонтитулы с титульного листа

\begin{center}
	\textbf{
		\university \\
		\institute \\
		\school\\
	}
    
	\vspace{2ex}

    \begin{flushright}
        \begin{minipage}{0.36\linewidth}
            Работа допущена к защите \\
            Директор ВШТМиМФ, \\
            д.ф.-м.н., чл.-корр. РАН \\
            \underline{\hspace{3cm}} А.М. Кривцов \\
            «\underline{\hspace{1cm}}» \underline{\hspace{3cm}} \publishYear \ г.
        \end{minipage}
    \end{flushright}
    
	\vspace{2ex}

	\textbf{\thesisTitle} \\
	{\large\textbf{\title}} \\
	{\large по направлению подготовки} \\
    {\large \studyField} \\
    {\large Направленность (профиль)} \\
    {\large \profile} \\
	
	\vspace{2ex}

	\begin{tabular}{p{0.4\textwidth} p{0.25\textwidth} p{0.2\textwidth}}
		\begin{flushleft}Выполнил \\ студент гр. \group\end{flushleft} & & \begin{flushleft}\author\end{flushleft} \\
		\begin{flushleft}Научный руководитель \\ \supervisorTitle  \end{flushleft}& & \begin{flushleft}\supervisor\end{flushleft}\\ 
	\end{tabular}
\end{center}

\vfill

\begin{center}
	\city \\
	\publishYear
\end{center}

\newpage
	% Титульник
    % \includepdf[pages=-,  scale=1]{pdfs/01_Титульник.pdf}
	% Индивидуальный план
    % \includepdf[pages=-,  scale=1]{pdfs/02_ИП.pdf}

	% 2. Задание на ВКР

	% 3. Реферат
	\newpage
	\section*{РЕФЕРАТ}
		\setcounter{section}{0}
		\setcounter{subsection}{0}
		\setcounter{figure}{0}
		\setcounter{table}{0}
		\subfile{sections/30_Реферат}

	% 4. Содержание
	\newpage
	\renewcommand{\contentsname}{\normalfont\centering\bfseries СОДЕРЖАНИЕ}
	\singlespacing  % Change spacing only for TOC
	\tableofcontents
	\onehalfspacing
	\thispagestyle{empty}

	% 5. Введение
	\newpage
	\section*{ВВЕДЕНИЕ}
		\addcontentsline{toc}{section}{Введение}
		\addtocontents{toc}{\protect\thispagestyle{empty}}  % To delete page number from toc
		\setcounter{section}{0}
		\setcounter{subsection}{0}
		\setcounter{figure}{0}
		\setcounter{table}{0}
		\subfile{sections/50_Введение}

	% 6. Основная часть
	% !TEX root = ../main.tex

\newpage
\section*{ГЛАВА 1. НАЗВАНИЕ ГЛАВЫ}
\addcontentsline{toc}{section}{Глава 1. Название главы}
\setcounter{section}{1} 
\setcounter{subsection}{0}
\setcounter{figure}{0}
\setcounter{table}{0}
\subfile{sections/61_Глава_1}

\newpage
\section*{ГЛАВА 2. НАЗВАНИЕ ГЛАВЫ}
\addcontentsline{toc}{section}{Глава 2. Название главы}
\setcounter{section}{2} 
\setcounter{subsection}{0}
\setcounter{figure}{0}
\setcounter{table}{0}
\subfile{sections/62_Глава_2}

\newpage
\section*{ГЛАВА 3. НАЗВАНИЕ ГЛАВЫ}
\addcontentsline{toc}{section}{Глава 3. Название главы}
\setcounter{section}{3} 
\setcounter{subsection}{0}
\setcounter{figure}{0}
\setcounter{table}{0}
\subfile{sections/63_Глава_3}


	% 7. Заключение
	\newpage
	\section*{ЗАКЛЮЧЕНИЕ}
		\addcontentsline{toc}{section}{Заключение}
		\setcounter{section}{0}
		\setcounter{subsection}{0}
		\setcounter{figure}{0}
		\setcounter{table}{0}
		\subfile{sections/70_Заключение}

	% 8. Список использованных источников
	\newpage
\printbibliography[title=СПИСОК ИСПОЛЬЗОВАННЫХ ИСТОЧНИКОВ]
\addcontentsline{toc}{section}{Список использованных источников}


	% 9. Приложения
	% !TEX root = main.tex

%\vspace*{40ex}
%\appendix
%\section*{ПРИЛОЖЕНИЯ}
\addcontentsline{toc}{section}{Приложения}
\renewcommand{\thesubsection}{\Alph{subsection}.}  % Буквы в  subsections
\renewcommand\thetable{\arabic{subsection}.\arabic{table}} 
\setcounter{section}{0} 
\setcounter{subsection}{0}
\setcounter{figure}{0}
\setcounter{table}{0}

% Change header height
%\setlength{\headheight}{3\baselineskip}
%\setlength{\headsep}{12pt}
%\fancyhead[L]{\rule{30pt}{2pt}}}

\addcontentsline{toc}{subsection}{Приложение 1. Название приложения}
\setcounter{subsection}{1}
\includepdf[pages=-,  scale=1]{pdfs/pdf1.pdf}

\newgeometry{left=20mm, top=20mm, right=20mm, bottom=20mm}



\newpage
%\rhead{\fontsize{14pt}{16pt}\selectfont\thepage\linebreak Приложение 2}
\rhead{\fontsize{14pt}{16pt}\selectfont\thepage}
Приложение 2\raggedleft
\addcontentsline{toc}{subsection}{Приложение 2. Название приложения}
\setcounter{subsection}{2}

\begin{center}
	\textbf{Название приложения}
\end{center}

\justifying
Текст приложения...


\newpage
%\rhead{\fontsize{14pt}{16pt}\selectfont\thepage\linebreak Приложение 3}
\rhead{\fontsize{14pt}{16pt}\selectfont\thepage}
Приложение 3 \raggedleft
\addcontentsline{toc}{subsection}{Приложение 3. Название приложения}
\setcounter{subsection}{3}

\begin{center}
	\textbf{Название приложения}
\end{center}

\justifying
Текст приложения...








\end{document}


